\sisetup{output-decimal-marker = {,}}

\section{Introdução}
% TODO

\begin{table}[htbp]
    \centering
    \renewcommand{\thetable}{C}\renewcommand{\theHtable}{C}
    \renewcommand{\arraystretch}{1.35}\setlength{\tabcolsep}{5pt}
    \begin{tabular}{l@{\hspace{1.6em}}rrr}
        \toprule
         & inversor & avanço & RC \\
        \midrule
        $R$ & -- & -- & 10 k$\Omega$ \\
        $R_{1}$ & -- & 12 k$\Omega$ & 15 k$\Omega$ \\
        $C_{1}$ & -- & 1 $\mu$F & 47 nF \\
        $R_{2}$ & -- & 10 k$\Omega$ & 390 k$\Omega$ \\
        $C_{2}$ & -- & 100 nF & 1 nF \\
        $R_{3}$ & 10 k$\Omega$ & 50 k$\Omega$ & -- \\
        $R_{4}$ & 10 k$\Omega$ & 10 k$\Omega$ & -- \\
        $R_{5}$ & -- & -- & 10 k$\Omega$ \\
        $R_{6}$ & -- & -- & 10 k$\Omega$ \\
        $R_{c}$ & -- & -- & 1 k$\Omega$ \\
        $C_{c}$ & -- & -- & 1 $\mu$F \\
        \bottomrule
    \end{tabular}
    \caption{Componentes nominais de cada montagem (valores do roteiro).}
    \label{tab:componentes}
\end{table}



\section{Circuito inversor}

\subsection{Tabela 1}
% Monte o circuito inversor (R3 = R4 = 10 kΩ), aplique senoide de 500 Hz, 2 Vpp, e meça
% o Vpp da saída e o ângulo de fase.
% TODO

\begin{table}[htbp]
    \centering
    \setcounter{table}{0}% mesmo número da tabela do roteiro
    \renewcommand{\arraystretch}{1.35}\setlength{\tabcolsep}{5pt}
    \begin{tabular}{l@{\hspace{1.6em}}r@{\hspace{1.6em}}rr@{\hspace{1.6em}}rr@{\hspace{1.6em}}rr}
        \toprule
        & & \multicolumn{2}{c}{$V_{pp}$ da saída (V)} & \multicolumn{2}{c}{$|G|$} & \multicolumn{2}{c}{fase ($^\circ$)} \\
        \cmidrule(lr){3-4} \cmidrule(lr){5-6} \cmidrule(lr){7-8}
        caso & $V_{pp}$ entrada (V) & exp. & \textcolor{black!55}{prev.} & exp. & teór. & exp. & teór. \\
        \midrule
        500\,Hz & 2,006 & 2,009 & \textcolor{black!55}{2,006} & 1,002 & 1,000 & 179,9 & 180,0 \\
        \bottomrule
    \end{tabular}
    \par\smallskip{\footnotesize Previsto de $V_{pp}$ da saída: $|G|$ teórico $\times$ $V_{pp}$ de entrada medido em CH1. Fase de CH2 em relação a CH1, pelo ajuste de senoide.}
    \caption{Circuito inversor ($R_{3}$ = $R_{4}$ = 10 k$\Omega$, senoide de 500 Hz, 2\,V$_{pp}$).}
    \label{tab:tabela1}
\end{table}

% Cores validadas para daltonismo (Okabe-Ito); o traço também distingue as curvas na impressão P&B.
\definecolor{cor-medido}{HTML}{0072B2}
\definecolor{cor-teorico}{HTML}{D55E00}
\definecolor{cor-carga}{HTML}{009E73}
\begin{figure}[htbp]
    \centering
    \pgfplotslegendfromname{leg-sen1}\par\smallskip
    \begin{tikzpicture}
        \begin{groupplot}[
            group style={group size=1 by 1, horizontal sep=1.8cm, vertical sep=2.3cm},
            scale only axis, width=0.7\textwidth, height=0.2\textwidth,
            xlabel={tempo (ms)}, ylabel={tensão (V)},
            grid=major, grid style={gray!20},
            title style={font=\small, at={(0,1)}, anchor=south west},
            scaled ticks=false,
            tick label style={font=\footnotesize, /pgf/number format/.cd, fixed, use comma, 1000 sep={\,}},
            legend columns=3, legend style={font=\footnotesize, draw=none, /tikz/every even column/.append style={column sep=0.6em}},
            unbounded coords=jump,
        ]
            \nextgroupplot[title={500\,Hz \quad (captura \#02)}, xmin=0, xmax=4.999, legend to name=leg-sen1]
            \addplot[gray!80, thick] table[x=t, y=ch1, col sep=comma] {10/csv/senoide_t1_c1.1.csv}; \addlegendentry{entrada (CH1)}
            \addplot[cor-medido, line width=1pt] table[x=t, y=ch2, col sep=comma] {10/csv/senoide_t1_c1.1.csv}; \addlegendentry{saída medida (CH2)}
            \addplot[cor-teorico, dashed, line width=0.9pt] table[x=t, y=teo, col sep=comma] {10/csv/senoide_t1_c1.1.csv}; \addlegendentry{saída teórica}
        \end{groupplot}
    \end{tikzpicture}
    \caption{Casos da Tabela~\ref{tab:tabela1}: entrada (CH1), saída medida (CH2) e saída prevista
             pelo modelo para a mesma entrada.}
    \label{fig:senoide1}
\end{figure}



\section{Compensador avanço de fase em malha aberta}

\subsection{Função de transferência}
% TODO

\begin{equation}
    G_c(s) = \frac{R_2}{R_1}\,\frac{R_1 C_1\, s + 1}{R_2 C_2\, s + 1}
\end{equation}

\subsection{Tabela 2}
% Com C1 = 1 µF, C2 = 100 nF, R1 = 12 kΩ, R2 = 10 kΩ, R3 = 50 kΩ, R4 = 10 kΩ, aplique
% senoides de 2 Vpp em 20 Hz, 100 Hz, 300 Hz, 1 kHz, 3 kHz e 5 kHz e meça o Vpp da saída
% e o ângulo de fase.
% TODO

\begin{table}[htbp]
    \centering
    \setcounter{table}{1}% mesmo número da tabela do roteiro
    \renewcommand{\arraystretch}{1.35}\setlength{\tabcolsep}{5pt}
    \begin{tabular}{l@{\hspace{1.6em}}r@{\hspace{1.6em}}rr@{\hspace{1.6em}}rr@{\hspace{1.6em}}rr}
        \toprule
        & & \multicolumn{2}{c}{$V_{pp}$ da saída (V)} & \multicolumn{2}{c}{$|G|$} & \multicolumn{2}{c}{fase ($^\circ$)} \\
        \cmidrule(lr){3-4} \cmidrule(lr){5-6} \cmidrule(lr){7-8}
        caso & $V_{pp}$ entrada (V) & exp. & \textcolor{black!55}{prev.} & exp. & teór. & exp. & teór. \\
        \midrule
        20\,Hz & 1,003 & 0,296 & \textcolor{black!55}{0,300} & 0,295 & 0,299 & 48,3 & 49,3 \\
        100\,Hz & 0,998 & 1,049 & \textcolor{black!55}{1,072} & 1,051 & 1,073 & 53,5 & 50,3 \\
        300\,Hz & 0,994 & 1,964 & \textcolor{black!55}{1,758} & 1,976 & 1,768 & 31,2 & 25,4 \\
        1\,kHz & 0,945 & 2,257 & \textcolor{black!55}{1,867} & 2,388 & 1,975 & 9,0 & 8,3 \\
        3\,kHz & 0,713 & 1,629 & \textcolor{black!55}{1,425} & 2,283 & 1,997 & $-$2,7 & 2,8 \\
        5\,kHz & 0,529 & 1,153 & \textcolor{black!55}{1,058} & 2,179 & 1,999 & $-$8,0 & 1,7 \\
        \bottomrule
    \end{tabular}
    \par\smallskip{\footnotesize Previsto de $V_{pp}$ da saída: $|G|$ teórico $\times$ $V_{pp}$ de entrada medido em CH1. Fase de CH2 em relação a CH1, pelo ajuste de senoide.}
    \caption{Compensador avanço de fase ($C_{1}$ = 1 $\mu$F, $C_{2}$ = 100 nF, $R_{1}$ = 12 k$\Omega$, $R_{2}$ = 10 k$\Omega$, $R_{3}$ = 50 k$\Omega$, $R_{4}$ = 10 k$\Omega$, senoide de 2\,V$_{pp}$).}
    \label{tab:tabela2}
\end{table}

% Cores validadas para daltonismo (Okabe-Ito); o traço também distingue as curvas na impressão P&B.
\definecolor{cor-medido}{HTML}{0072B2}
\definecolor{cor-teorico}{HTML}{D55E00}
\definecolor{cor-carga}{HTML}{009E73}
\begin{figure}[htbp]
    \centering
    \pgfplotslegendfromname{leg-sen2}\par\smallskip
    \begin{tikzpicture}
        \begin{groupplot}[
            group style={group size=2 by 3, horizontal sep=1.8cm, vertical sep=2.3cm},
            scale only axis, width=0.4\textwidth, height=0.2\textwidth,
            xlabel={tempo (ms)}, ylabel={tensão (V)},
            grid=major, grid style={gray!20},
            title style={font=\small, at={(0,1)}, anchor=south west},
            scaled ticks=false,
            tick label style={font=\footnotesize, /pgf/number format/.cd, fixed, use comma, 1000 sep={\,}},
            legend columns=3, legend style={font=\footnotesize, draw=none, /tikz/every even column/.append style={column sep=0.6em}},
            unbounded coords=jump,
        ]
            \nextgroupplot[title={20\,Hz \quad (captura \#20)}, xmin=0, xmax=150, legend to name=leg-sen2]
            \addplot[gray!80, thick] table[x=t, y=ch1, col sep=comma] {10/csv/senoide_t2_c2.1.csv}; \addlegendentry{entrada (CH1)}
            \addplot[cor-medido, line width=1pt] table[x=t, y=ch2, col sep=comma] {10/csv/senoide_t2_c2.1.csv}; \addlegendentry{saída medida (CH2)}
            \addplot[cor-teorico, dashed, line width=0.9pt] table[x=t, y=teo, col sep=comma] {10/csv/senoide_t2_c2.1.csv}; \addlegendentry{saída teórica}
            \nextgroupplot[title={100\,Hz \quad (captura \#21)}, xmin=0, xmax=24.99]
            \addplot[gray!80, thick] table[x=t, y=ch1, col sep=comma] {10/csv/senoide_t2_c2.2.csv};
            \addplot[cor-medido, line width=1pt] table[x=t, y=ch2, col sep=comma] {10/csv/senoide_t2_c2.2.csv};
            \addplot[cor-teorico, dashed, line width=0.9pt] table[x=t, y=teo, col sep=comma] {10/csv/senoide_t2_c2.2.csv};
            \nextgroupplot[title={300\,Hz \quad (captura \#19)}, xmin=0, xmax=9.997]
            \addplot[gray!80, thick] table[x=t, y=ch1, col sep=comma] {10/csv/senoide_t2_c2.3.csv};
            \addplot[cor-medido, line width=1pt] table[x=t, y=ch2, col sep=comma] {10/csv/senoide_t2_c2.3.csv};
            \addplot[cor-teorico, dashed, line width=0.9pt] table[x=t, y=teo, col sep=comma] {10/csv/senoide_t2_c2.3.csv};
            \nextgroupplot[title={1\,kHz \quad (captura \#22)}, xmin=0, xmax=2.499]
            \addplot[gray!80, thick] table[x=t, y=ch1, col sep=comma] {10/csv/senoide_t2_c2.4.csv};
            \addplot[cor-medido, line width=1pt] table[x=t, y=ch2, col sep=comma] {10/csv/senoide_t2_c2.4.csv};
            \addplot[cor-teorico, dashed, line width=0.9pt] table[x=t, y=teo, col sep=comma] {10/csv/senoide_t2_c2.4.csv};
            \nextgroupplot[title={3\,kHz \quad (captura \#23)}, xmin=0, xmax=0.9997]
            \addplot[gray!80, thick] table[x=t, y=ch1, col sep=comma] {10/csv/senoide_t2_c2.5.csv};
            \addplot[cor-medido, line width=1pt] table[x=t, y=ch2, col sep=comma] {10/csv/senoide_t2_c2.5.csv};
            \addplot[cor-teorico, dashed, line width=0.9pt] table[x=t, y=teo, col sep=comma] {10/csv/senoide_t2_c2.5.csv};
            \nextgroupplot[title={5\,kHz \quad (captura \#24)}, xmin=0, xmax=0.4999]
            \addplot[gray!80, thick] table[x=t, y=ch1, col sep=comma] {10/csv/senoide_t2_c2.6.csv};
            \addplot[cor-medido, line width=1pt] table[x=t, y=ch2, col sep=comma] {10/csv/senoide_t2_c2.6.csv};
            \addplot[cor-teorico, dashed, line width=0.9pt] table[x=t, y=teo, col sep=comma] {10/csv/senoide_t2_c2.6.csv};
        \end{groupplot}
    \end{tikzpicture}
    \caption{Casos da Tabela~\ref{tab:tabela2}: entrada (CH1), saída medida (CH2) e saída prevista
             pelo modelo para a mesma entrada.}
    \label{fig:senoide2}
\end{figure}

% Cores validadas para daltonismo (Okabe-Ito); o traço também distingue as curvas na impressão P&B.
\definecolor{cor-medido}{HTML}{0072B2}
\definecolor{cor-teorico}{HTML}{D55E00}
\definecolor{cor-carga}{HTML}{009E73}
\begin{figure}[htbp]
    \centering
    \pgfplotslegendfromname{leg-bode2}\par\smallskip
    \begin{tikzpicture}
        \begin{groupplot}[
            group style={group size=1 by 2, vertical sep=0.6cm, xlabels at=edge bottom, xticklabels at=edge bottom},
            scale only axis, width=0.8\textwidth, height=0.22\textwidth, xmode=log,
            xlabel={frequência (Hz)},
            grid=major, grid style={gray!20},
            title style={font=\small, at={(0,1)}, anchor=south west},
            scaled ticks=false,
            tick label style={font=\footnotesize, /pgf/number format/.cd, fixed, use comma, 1000 sep={\,}},
            legend columns=3, legend style={font=\footnotesize, draw=none, /tikz/every even column/.append style={column sep=0.6em}},
            unbounded coords=jump,
        ]
            \nextgroupplot[ylabel={$|G|$ (dB)}, legend to name=leg-bode2]
            \addplot[cor-teorico, dashed, line width=0.9pt, mark=none] table[x=f, y=db, col sep=comma] {10/csv/bode_t2_teo.csv};
            \addlegendentry{teórico}
            \addplot[only marks, mark=square*, mark size=2.2pt, cor-medido] table[x=f, y=db, col sep=comma] {10/csv/bode_t2_exp.csv};
            \addlegendentry{medido}
            \nextgroupplot[ylabel={fase ($^\circ$)}]
            \addplot[cor-teorico, dashed, line width=0.9pt, mark=none] table[x=f, y=fase, col sep=comma] {10/csv/bode_t2_teo.csv};
            \addplot[only marks, mark=square*, mark size=2.2pt, cor-medido] table[x=f, y=fase, col sep=comma] {10/csv/bode_t2_exp.csv};
        \end{groupplot}
    \end{tikzpicture}
    \caption{Resposta em frequência do compensador (Tabela~\ref{tab:tabela2}): módulo e fase,
             modelo teórico e pontos medidos.}
    \label{fig:bode2}
\end{figure}



\section{Malha fechada com planta RC}

\subsection{Função de transferência}
% TODO

\subsection{Tabela 3}
% Com Cc = 1 µF, Rc = 1 kΩ, C1 = 47 nF, C2 = 1 nF, R1 = 15 kΩ, R2 = 390 kΩ, R = R6 = 10 kΩ,
% aplique onda quadrada de 50 Hz, 1 Vpp, e meça τ, Tr, Ts e Vc.
% TODO

\begin{table}[htbp]
    \centering
    \setcounter{table}{2}% mesmo número da tabela do roteiro
    \renewcommand{\arraystretch}{1.35}\setlength{\tabcolsep}{5pt}
    \begin{tabular}{l@{\hspace{1.6em}}rrr@{\hspace{1.6em}}r}
        \toprule
        métrica & exp. & teór. & \textcolor{black!55}{c/ carga} & desvio \\
        \midrule
        $\tau$ ($\mu$s) & 61 & 21 & \textcolor{black!55}{21} & 197,2\,\% \\
        $T_r$ ($\mu$s) & 66 & 46 & \textcolor{black!55}{46} & 43,9\,\% \\
        $T_s$ ($\mu$s) & 114 & 87 & \textcolor{black!55}{85} & 30,6\,\% \\
        $V_c$ (V) & 0,479 & 0,481 & \textcolor{black!55}{0,480} & $-$0,5\,\% \\
        \bottomrule
    \end{tabular}
    \par\smallskip{\footnotesize $T_r$ de 10\,\% a 90\,\%, $T_s$ na faixa de $\pm 2$\,\%, $\tau$ até 63,2\,\%; a saída medida foi suavizada por média móvel de 5,7\,$\mu$s.}
    \caption{RC com compensador avanço de fase ($C_{c}$ = 1 $\mu$F, $R_{c}$ = 1 k$\Omega$, $C_{1}$ = 47 nF, $C_{2}$ = 1 nF, $R_{1}$ = 15 k$\Omega$, $R_{2}$ = 390 k$\Omega$, $R$ = $R_6$ = 10 k$\Omega$).}
    \label{tab:tabela3}
\end{table}

% Cores validadas para daltonismo (Okabe-Ito); o traço também distingue as curvas na impressão P&B.
\definecolor{cor-medido}{HTML}{0072B2}
\definecolor{cor-teorico}{HTML}{D55E00}
\definecolor{cor-carga}{HTML}{009E73}
\begin{figure}[htbp]
    \centering
    \pgfplotslegendfromname{leg-deg3}\par\smallskip
    \begin{tikzpicture}
        \begin{groupplot}[
            group style={group size=1 by 1, horizontal sep=1.8cm, vertical sep=1.6cm},
            scale only axis, width=0.7\textwidth, height=0.24\textwidth,
            xlabel={tempo após a borda de subida (ms)}, ylabel={tensão (V)},
            grid=major, grid style={gray!20},
            title style={font=\small, at={(0,1)}, anchor=south west},
            scaled ticks=false,
            tick label style={font=\footnotesize, /pgf/number format/.cd, fixed, use comma, 1000 sep={\,}},
            legend columns=3, legend style={font=\footnotesize, draw=none, /tikz/every even column/.append style={column sep=0.6em}},
            unbounded coords=jump, legend columns=4,
        ]
            \nextgroupplot[xmin=-0.017, xmax=0.205, legend to name=leg-deg3]
            \addplot[gray!70, thick, const plot] table[x=t, y=e1, col sep=comma] {10/csv/degrau_t3_c3.1.csv}; \addlegendentry{entrada $e_1$}
            \addplot[cor-medido, line width=1.1pt] table[x=t, y=exp, col sep=comma] {10/csv/degrau_t3_c3.1.csv}; \addlegendentry{medido ($-V_c$)}
            \addplot[cor-teorico, dashed, line width=0.9pt] table[x=t, y=teo, col sep=comma] {10/csv/degrau_t3_c3.1.csv}; \addlegendentry{teórico}
            \addplot[cor-carga, densely dotted, line width=1.1pt] table[x=t, y=carga, col sep=comma] {10/csv/degrau_t3_c3.1.csv}; \addlegendentry{c/ carga do inversor}
        \end{groupplot}
    \end{tikzpicture}
    \caption{Resposta ao degrau da Tabela~\ref{tab:tabela3}: entrada, saída medida e saídas dos dois
             modelos, todas com o sinal trocado ($-V_c$), já que a malha inverte a saída em relação à entrada.}
    \label{fig:degrau3}
\end{figure}


\subsection{Tabela 4}
% Mesmo circuito, senoide de 100 Hz, 1 Vpp: Vpp da saída e ângulo de fase.
% TODO

\begin{table}[htbp]
    \centering
    \setcounter{table}{3}% mesmo número da tabela do roteiro
    \renewcommand{\arraystretch}{1.35}\setlength{\tabcolsep}{5pt}
    \begin{tabular}{l@{\hspace{1.6em}}r@{\hspace{1.6em}}rr@{\hspace{1.6em}}rr@{\hspace{1.6em}}rr}
        \toprule
        & & \multicolumn{2}{c}{$V_{pp}$ da saída (V)} & \multicolumn{2}{c}{$|G|$} & \multicolumn{2}{c}{fase ($^\circ$)} \\
        \cmidrule(lr){3-4} \cmidrule(lr){5-6} \cmidrule(lr){7-8}
        caso & $V_{pp}$ entrada (V) & exp. & \textcolor{black!55}{prev.} & exp. & teór. & exp. & teór. \\
        \midrule
        RC, 100\,Hz & 1,007 & 0,967 & \textcolor{black!55}{0,969} & 0,960 & 0,962 & 179,2 & 179,1 \\
        \bottomrule
    \end{tabular}
    \par\smallskip{\footnotesize Previsto de $V_{pp}$ da saída: $|G|$ teórico $\times$ $V_{pp}$ de entrada medido em CH1. Fase de CH2 em relação a CH1, pelo ajuste de senoide.}
    \caption{Avanço de fase, circuito RC (senoide de 100 Hz, 1 V$_{pp}$).}
    \label{tab:tabela4}
\end{table}

% Cores validadas para daltonismo (Okabe-Ito); o traço também distingue as curvas na impressão P&B.
\definecolor{cor-medido}{HTML}{0072B2}
\definecolor{cor-teorico}{HTML}{D55E00}
\definecolor{cor-carga}{HTML}{009E73}
\begin{figure}[htbp]
    \centering
    \pgfplotslegendfromname{leg-sen4}\par\smallskip
    \begin{tikzpicture}
        \begin{groupplot}[
            group style={group size=1 by 1, horizontal sep=1.8cm, vertical sep=2.3cm},
            scale only axis, width=0.7\textwidth, height=0.2\textwidth,
            xlabel={tempo (ms)}, ylabel={tensão (V)},
            grid=major, grid style={gray!20},
            title style={font=\small, at={(0,1)}, anchor=south west},
            scaled ticks=false,
            tick label style={font=\footnotesize, /pgf/number format/.cd, fixed, use comma, 1000 sep={\,}},
            legend columns=3, legend style={font=\footnotesize, draw=none, /tikz/every even column/.append style={column sep=0.6em}},
            unbounded coords=jump,
        ]
            \nextgroupplot[title={RC, 100\,Hz \quad (captura \#41)}, xmin=0, xmax=24.98, legend to name=leg-sen4]
            \addplot[gray!80, thick] table[x=t, y=ch1, col sep=comma] {10/csv/senoide_t4_c4.1.csv}; \addlegendentry{entrada (CH1)}
            \addplot[cor-medido, line width=1pt] table[x=t, y=ch2, col sep=comma] {10/csv/senoide_t4_c4.1.csv}; \addlegendentry{saída medida (CH2)}
            \addplot[cor-teorico, dashed, line width=0.9pt] table[x=t, y=teo, col sep=comma] {10/csv/senoide_t4_c4.1.csv}; \addlegendentry{saída teórica}
        \end{groupplot}
    \end{tikzpicture}
    \caption{Casos da Tabela~\ref{tab:tabela4}: entrada (CH1), saída medida (CH2) e saída prevista
             pelo modelo para a mesma entrada.}
    \label{fig:senoide4}
\end{figure}



\clearpage
\section{Malha fechada com planta RLC}
% Tabelas 5, 6 e 7: sem captura ainda. Depois de medir, rode o exportador de novo e acrescente
% \input{10/tab/tabela5.tex}, \input{10/tex/plot_degrau5.tex}, \input{10/tab/tabela6.tex},
% \input{10/tab/tabela7.tex} e \input{10/tex/plot_senoide7.tex}.
% TODO


\section{Conclusão}
% TODO
